\documentclass[11pt,a4paper]{article}

\usepackage[margin=2.5cm]{geometry}
\usepackage[T1]{fontenc}
\usepackage[utf8]{inputenc}
\usepackage{lmodern}
\usepackage{amsmath}
\usepackage{graphicx}
\usepackage{booktabs}      % nice tables: \toprule, \midrule, \bottomrule
\usepackage{subcaption}    % figures side by side
\usepackage[hidelinks]{hyperref}

\title{MALIS Project 1: Reinforcement Learning and Multi-Armed Bandits}
\author{Carlos Lugo Torres \and Axel Castro Lara \and Emilio Gonzalez Acosta}
\date{EURECOM -- November 2026}

\begin{document}
\maketitle

% ---------------------------------------------------------------
\section{Introduction}               % ~75 words
% The game (5 spots, noisy catches, 3 seasons) and the goal of the report.
\section{Introduction}

In \emph{Fishing Season}, a player chooses one of five fishing spots (the arms)
each day and observes the catch (the reward). Each spot has a hidden mean catch
between 2 and 10\,kg, and daily catches vary around it with a standard deviation
of 2.5\,kg. This is a multi-armed bandit: every day, the player must trade off
exploiting the spot that looks best against exploring the others. The game has
three seasons: 30 days, 10 days, and 30 days in which the means change on day 16.

We analyze our team's play (Section~\ref{sec:human}), propose a measure that
separates decisions from luck (Section~\ref{sec:measure}), implement our best
human strategy and two methods from Sutton and Barto~\cite{sutton}
(Section~\ref{sec:agents}), evaluate them on 1000 random bays per season
(Section~\ref{sec:results}), and extend the problem to a contextual bandit
(Section~\ref{sec:contextual}).

% ---------------------------------------------------------------
\section{Human play}
\label{sec:human}

\paragraph{Axel.} Every season I planned to use at most half of the days for
reconnaissance, to tell ``good spots'' from ``bad spots''. Then I stayed as long
as possible in good spots and, whenever I had a bad day, immediately changed to
another good spot. My log sometimes shows something different, because I also
tested how random the catches were, looking for a pattern.

\paragraph{Emilio.} In Seasons 1 and 3 I tried every spot once (days 1--5) and
then stayed at the best one, fishing Lighthouse Reef twice more in Season 1 to
confirm it was not an outlier. In Season 2, exploring all five spots would have
cost half the season, so I tried one spot at a time and stayed at the first one
with a high catch. In Season 3, after 0.5\,kg at Old Pier on day 21, I waited
five more days before switching to Rocky Point, my second best spot.

\paragraph{Carlos.} I did not follow an explicit plan; looking at my choices, my
strategy was to keep moving between the spots with good catches instead of
committing to one, and I almost never fished the same spot two days in a row.
In Seasons 2 and 3 I first tried every spot once; in Season 1 I mostly
alternated between Deep Channel and Seagrass Bed. After the change in Season 3,
I added Lighthouse Reef to the rotation once it gave 15.9\,kg on day 19.

\paragraph{Rules.} We abstract these behaviours into three rules that can be
applied to any log:
\begin{itemize}
  \item \textbf{R1, commit with patience $k$} (Emilio, Seasons 1 and 3; Axel):
  fish each spot once, then commit to the spot with the highest average. A catch
  is \emph{bad} if it is below that spot's own average; after $k$ bad catches in
  a row, move to the other spot with the highest average.
  \item \textbf{R2, sequential search with threshold $T$} (Emilio, Season 2):
  try spots one at a time until a catch reaches $T$, then stay there.
  \item \textbf{R3, rotation} (Carlos): never fish the same spot two days in a
  row, choosing among the spots with the best catches so far.
\end{itemize}

We check each rule by replaying the logs and counting the days (from day 6) on
which it predicts the player's choice. R1 with $k=6$ predicts all 25 of Emilio's
choices in Seasons 1 and 3, and R1 with $k=1$ predicts 16 and 19 of Axel's 25.
R2 with $T=8$\,kg reproduces Emilio's Season 2, and R3 never repeating a spot
matches 63 of Carlos's 67 day-to-day choices. As Emilio obtained the best score
in Seasons 1 and 2 (Section~\ref{sec:measure}), R1 is the strategy we implement
as our human agent (Section~\ref{sec:agents}).
% ---------------------------------------------------------------
\section{Separating decisions from luck}   % ~275 words  (Task 3)
\label{sec:measure}
\section{Separating decisions from luck}
\label{sec:measure}

The total catch mixes decisions with two kinds of luck: the daily noise of the
catches, and the bay itself, since some bays contain only good spots. We remove
the first by scoring the true mean of the chosen spot instead of the catch, and
the second by rescaling each day between the worst and the best spot:
\begin{equation}
  S = \frac{1}{T} \sum_{t=1}^{T}
      \frac{\mu_t(a_t) - \min_a \mu_t(a)}{\max_a \mu_t(a) - \min_a \mu_t(a)},
\end{equation}
where $a_t$ is the spot chosen on day $t$ and $\mu_t$ are the true means in
force that day (they change on day 16 of Season 3). $S=1$ means the best spot
every day and $S=0$ the worst; a wrong choice costs in proportion to its
distance from the best spot, and averaging over days makes seasons of different
lengths comparable. If all means are equal, the day scores 1.

To show that $S$ separates decisions from luck, we play every agent of
Section~\ref{sec:agents} on the same 1000 random bays per season
(Table~\ref{tab:measure}). The total catch follows the bay: for UCB in Season 1,
its correlation with the expected catch of a random player on the same bay is
0.62, against 0.06 for $S$. $S$ also separates good from bad decisions more
clearly: the gap between UCB and the random agent grows from 1.6 to 3.2
standard deviations, and UCB beats the random agent on 99\% of the bays
according to $S$, against 96\% according to the catch. One limitation remains:
the score of a random player still depends on the bay (correlation 0.42 in
Season 1), because it depends on how the means are spread between the worst and
the best spot.

\begin{table}[t]
\centering
\small
\begin{tabular}{lcccccc}
\toprule
 & \multicolumn{2}{c}{Corr.\ with bay quality (UCB)}
 & \multicolumn{2}{c}{UCB beats Random}
 & \multicolumn{2}{c}{Gap UCB -- Random (s.d.)} \\
\cmidrule(lr){2-3}\cmidrule(lr){4-5}\cmidrule(lr){6-7}
Season & kg & $S$ & kg & $S$ & kg & $S$ \\
\midrule
1 & 0.62 & 0.06 & 96\% & 99\% & 1.58 & 3.23 \\
2 & 0.68 & 0.31 & 75\% & 88\% & 0.67 & 1.34 \\
3 & 0.62 & 0.20 & 86\% & 88\% & 1.03 & 1.56 \\
\bottomrule
\end{tabular}
\caption{Total catch (kg) versus our score $S$ on 1000 random bays per season.
Bay quality is the expected total catch of a random player on that bay.}
\label{tab:measure}
\end{table}

\begin{table}[t]
\centering
\small
\begin{tabular}{lccccccc}
\toprule
 & \multicolumn{3}{c}{$S$ per season} & \multicolumn{2}{c}{Season 3}
 & New best spot \\
\cmidrule(lr){2-4}\cmidrule(lr){5-6}
 & 1 & 2 & 3 & days 1--15 & days 16--30 & first fished \\
\midrule
Axel   & 0.62 & 0.37 & 0.73 & 0.62 & 0.84 & never \\
Emilio & 0.92 & 0.88 & 0.55 & 0.78 & 0.33 & day 26 \\
Carlos & 0.59 & 0.70 & 0.43 & 0.41 & 0.45 & day 19 \\
\midrule
UCB (mean, 1000 bays) & 0.85 & 0.68 & 0.70 & & & \\
\bottomrule
\end{tabular}
\caption{Scores of the team members on their own bays.}
\label{tab:team}
\end{table}

\paragraph{Team.} Table~\ref{tab:team} gives our scores. In the short Season 2,
exploring every spot costs half of the season: Axel explored the whole season and
dropped to 0.37, while Emilio's sequential search kept 0.88. The agents show the
same effect (UCB falls from 0.85 to 0.68), since less time is left to exploit
what exploration found. In Season 3, all five means are redrawn on day 16, and
the best spot changed in all our bays. Axel, who left a spot after one bad day,
adapted fastest (0.84 after the change). Emilio stayed ten more days at Old
Pier, which had become the worst spot, and fell from 0.78 to 0.33. Carlos found
the new best spot on day 19, but his rotation kept him there only 4 of the last
15 days, so his score barely changed. Committing pays when the bay is stable but
slows the reaction to a change: this is the trade-off controlled by $k$ in our
human agent.

% ---------------------------------------------------------------
\section{Agents}
\label{sec:agents}

All agents are subclasses of \texttt{Agent} in \texttt{agents.py}: they only
observe the catch at the spot they chose, and \texttt{reset()} clears their
memory at the start of each season. All three begin by fishing each spot once,
in order, so that every spot has an estimate; ties are broken towards the lowest
spot number.

\paragraph{Human rule (R1).} \texttt{HumanRuleAgent} implements rule R1 of
Section~\ref{sec:human}. It stores the sum and number of catches of each spot,
the spot it is committed to, and the number of bad catches in a row at that
spot. After the first five days it commits to the spot with the highest average,
and it switches to the best other spot after $k$ bad catches in a row. The
patience $k$ trades robustness to bad luck against speed of reaction to a
change. We use $k=6$, the value that reproduces Emilio's play.

\paragraph{$\varepsilon$-greedy with constant step size.}
\texttt{EpsilonGreedyAlphaAgent} keeps an estimate $Q(a)$ and a count $N(a)$ for
each spot~\cite[Sec.~2.2--2.5]{sutton}. Each day it fishes a random spot with
probability $\varepsilon$, and the spot with the highest $Q$ otherwise. After a
catch $R$ at spot $a$, it updates
\begin{equation}
  Q(a) \leftarrow Q(a) + \alpha \bigl( R - Q(a) \bigr).
\end{equation}
With a constant $\alpha$, recent catches weigh more than old ones, which lets
the agent follow the change of Season 3; with $\alpha = 1/N(a)$
(\texttt{alpha=None}), $Q(a)$ is the sample average. The first catch at a spot
sets $Q(a)$ directly, which removes the bias of the initial
value~\cite[Sec.~2.6]{sutton}. We use $\varepsilon = 0.1$ and $\alpha = 0.5$.

\paragraph{Upper confidence bound (UCB).} \texttt{UCBAgent} keeps the sample
average $Q(a)$, the count $N(a)$ and the day $t$~\cite[Sec.~2.7]{sutton}. Each
day it fishes
\begin{equation}
  A_t = \arg\max_a \Bigl[ Q(a) + c \sqrt{\frac{\ln t}{N(a)}} \Bigr],
\end{equation}
so it explores the spots whose estimate is still uncertain (few visits) instead
of exploring at random. The parameter $c$ sets how much it explores. Since
catches vary with a standard deviation of 2.5\,kg, its useful scale differs
from the book's examples; we use $c = 1$.

% ---------------------------------------------------------------
\section{Results}                     % ~400 words  (Task 5)
\label{sec:results}
% 1000 random bays per season: distributions (figure + table).
% Agents vs. team members on their own bays.
% Best agent for Season 2 and Season 3, and why.
% Parameter study: epsilon, per season.

% ---------------------------------------------------------------
\section{Contextual bandit}           % ~200 words  (Task 6, optional)
\label{sec:contextual}
% Linear regression agent, comparison with UCB,
% failure case (bay 972) and the fix (warm-up m), with its trade-off.

% ---------------------------------------------------------------
\section{Conclusion}                  % ~50 words

% ===============================================================
% Not counted in the word limit
% ===============================================================
\section*{Contributions}
% What each member did.

\section*{Use of AI assistants}
% Required: which tool, what it was used for, which code it wrote.

\begin{thebibliography}{9}
\bibitem{sutton} R.~S. Sutton and A.~G. Barto,
  \emph{Reinforcement Learning: An Introduction}, 2nd ed., MIT Press, 2018.
\end{thebibliography}

\end{document}